\def\velkostpisma{10}
\def\velkostriadku{12.5}
\input makra.tex % nacitanie Ivanom pripravenych nastaveni a prikazov
\hyphenation{star-šov-stvo} % rozdelenie slov na konci riadku, treba tu uviest slova, ktore sam nepozna

\spravodaj{5}{2026}


\clanok{Verš na mesiac}
V~máji sa budeme učiť nový verš, ktorý dostali pre rok~2026 naše sestry:

„Vy ste však vyvolený rod, kráľovské kňazstvo, svätý národ, ľud určený na vlastníctvo, aby ste oznámili veľké skutky, toho, čo vás povolal z~temnoty do svojho predivného svetla.“ (1Pt~2,9)


\clanok{{\it ČO NÁS ČAKÁ V~MÁJI}}


\clanok{Večer chvál}
Chválospevkový tím nášho zboru organizuje 8.~5. Večer chvál. Príďte o~18.00~hod. do našej modlitebne chváliť nášho Pána.


\clanok{Deň otvorených dverí v~BJB Palisády}
V~sobotu 9.~5. od~14.00 do~18.00~hod. sa uskutoční Deň otvorených dverí na Súľovskej~2. Bude to dobrá príležitosť spolu stráviť čas a osloviť naše neveriace okolie prirodzeným spôsobom.


\clanok{Spoločný obed v~SED}
V~máji máme naplánovaný tretí náš spoločný obed. Stretneme sa v~nedeľu 10.~5. o~12.00~hod. v~priestoroch SED na Partizánskej. Prihlasovať sa nemusíte. Prineste si svoj obed a strávte tento čas pri spoločnom stolovaní a rozhovoroch.


\clanok{Rozhovory na Palisádach}
V~sobotu 16.~5.~2026 o~9.30~hod. na Palisádach prebehne ďalší z~„Rozhovorov na Palisádach“. Rečníčkou bude Darina Hlubocky (Malá) s~témou „Single, a predsa nie slobodná. Zadaný, a predsa single.“ Pozvánka s~viac informáciami je aj vo vašich e-mailových schránkach. Ak vás táto téma zaujíma, ste srdečne vítaní. Pozvite aj svojich priateľov.


\clanok{Sestry}
Sestry sa stretnú tento mesiac v~stredu 20.~5. o~17.30~hod. na Partizánskej~2 v~SED. Hosťom bude s.~Beáta Dobová s~témou: „Ticho pri studni“(…príbeh ženy, ktorá už nechcela byť smädná).


\clanok{Chvojnica na Letnice}
Víkend 23.~--~24.~5., na sviatok Zoslania Ducha Svätého, organizujeme tradičný výlet na Chvojnicu. Prihlasovanie sme spustili cez tabuľku v~e-maily. Ak máte vôľu, čas a chuť, pripojte sa v~sobotu k~jarnému upratovaniu chalupy a okolia. Nedeľu oslávime bohoslužbou v~kostolíku a sadneme si k~spoločnému obedu v~stodole.
Nech sa páči, pridajte sa.


\clanok{{\it OHLIADNUTIE SA ZA APRÍLOM}}


\clanok{Správy zo staršovstva}
V~minulom mesiaci sa staršovstvo stretlo dvakrát. Na prvom stretnutí sme sa rozprávali aj so sestrou Slávkou~V. o~potrebách a budúcnosti služby varenia polievky pre bezdomovcov. Sme veľmi vďační za bratov a sestry, ktorí sa pravidelne zapájajú do tejto služby. Tešíme sa, že sestra Slávka túto službu za náš zbor verne a ochotne zastrešuje. Upozornila nás však na potrebu obnovenia duchovnej služby počas výdaja. Dohodli sme, že brat kazateľ Janko bude túto potrebu tlmočiť na stretnutí pastorov v~rámci KvM.

Ďalším hosťom bol brat kazateľ Tomáš Valchář, s~ktorým sme sa rozprávali o~možnostiach spolupráce s~našou misijnou stanicou Connect ako aj o~spoločnom prevzatí správcovstva priestorov na Súľovskej~2 od~1.~4.~2026.

Ďalšími témami tohto stretnutia bolo aj zhodnotenie VZČZ (výročného zborového členského zhromaždenia) a zároveň prvého kola volieb kazateľa nášho zboru, podpora denného letného tábora pre deti, rekonštrukcia fasády nášho kostola, termíny KDZ (konferencií delegátov zborov).

Druhé stretnutie bolo po ukončení volieb kazateľa zboru a sme vďační jednak za zvolenie brata Dávida M.~Chuchúta za kazateľa nášho zboru ako aj za celkový pokojný priebeh týchto volieb.

Ďalšie dôležité témy stretnutia boli praktické kroky vyplývajúce zo spoločného prevzatia správy priestorov na Súľovskej -- návrh zmluvy ako aj kontaktná osoba za náš zbor -- br.~Dávid M.~Chuchút, rekonštrukcia fasády (s~plánovaným začiatkom prác od~1.~4.) a dohodnutie tém a cieľov zborovej víkendovky. Aj na tomto stretnutí sme privítali hosťa br.~Ľuboša Kešjara, s~ktorým sme sa rozprávali o~službe technikov nášho zboru ako aj o~tom, akým spôsobom môžeme zlepšiť vzájomnú spoluprácu v~kontexte nedeľných bohoslužieb ako aj pri iných akciách, ktoré vyžadujú zapojenie bratov z~techniky.

Všetkým, ktorí na nás v~modlitbách myslíte, ďakujeme a prosíme, aby ste v~tom aj naďalej pokračovali, potrebujeme múdrosť od Pána Boha a vedenie Duchom Svätým.
\autor{Radislav Nemec}


\clanok{Naše veľkonočné koncerty}
Všimli sme si, že piesne nášho spevokolu priťahujú do našej modlitebne veľa ľudí. A~hlavne neznámych a cudzích. Keďže nechodíme evanjelizovať do ulíc mesta, Pán Boh nám posiela ľudí priamo k~nám. Tak sa stalo pomaly tradíciou, že náš spevokol pripravuje koncerty s~vianočnou a veľkonočnou tematikou.

Keď sa pred pár rokmi stalo, že sa naša sála na Palisádach zaplnila tak, že asi päťdesiat ľudí sa nedostalo ani dnu, rozhodli sme urobiť dva rovnaké koncerty. Jeden v~sobotu a ten istý ešte raz hneď na druhý deň. Celkom rýchlo sme si na to zvykli. Pred rokom nás potom oslovilo vedenie ECAV z~Petržalky, aby sme ten druhý koncert mali u~nich. Veľmi sme sa tomu potešili.

Niekto si môže pomyslieť, že to pomaly môže vojsť až do stereotypu, ale Pán Boh to nedovolil. Tento rok nás dokonale preskúšal. My sme sa poctivo pripravovali a všetko si naplánovali ako pred každým koncertom. Ale keď sa blížil ich termín, začali mi chodiť správy o~zdravotných problémoch kľúčových hudobníkov a spevákov. Postihlo to sólistu i polovicu našich basistov. Prišlo to až tak ďaleko, že v~predvečer prvého koncertu som telefonoval s~dcérou Diankou, že nevidím iné východisko, ako oba koncerty zrušiť. Ona je kreatívna a máva dobré nápady. Teraz len stále opakovala, že ich zrušiť nemôžeme. Cez noc som vyradil zo zoznamu dve piesne a jednu upravil tak, aby sme to mohli spievať v~takom stave, v~akom sme. A~čakal ako bude ráno.

Správa, v~akej sme situácii, sa rýchlo rozšírila medzi spevákov. Keď som neskôr komunikoval o~pocitoch niektorých z~nich, jedna sestra mi napísala:

„No ja Ti poviem Slávko, že keď mi Ingrid začala na zastávke MHD v~sobotu pred treťou zvestovať všetky ‚Jóbove zvesti‘, až mi prišlo zle. Hlavou mi bežalo: ‚PANE, tak som sa modlila, bolo to málo? Prečooo? Tak úprimne som prosila...‘ Ale potom som to v~autobuse predýchala a keď som Ťa uvidela už v~kostole, vedela som, že sa na Toho hore a na Teba sa môžem spoľahnúť -- dopadne to dobre a bude to na oslavu Božiu. A~tak aj bolo. Uvedomujem si, akí sme krehkí!“

Všetci sme vďační, že tie koncerty sme nezrušili. Chceli sme nimi prinášať požehnanie a zvláštnou premenou prišlo požehnanie na nás. Veď aj pre nás platia slová 1.~Žalmu: „Blahoslavený ten, kto nechodí podľa cesty bezbožných…a všetko čo robí, podarí sa mu.“
\autor{Slávo Kráľ, dirigent spevokolu}


\clanok{Sedem slov}
Počula som ten príbeh už mnohokrát. S~trochou odvahy by som ho vedela prerozprávať aj ja. A~tentokrát bol predsa úplne iný.
Jednoduchý.

Bez pridaných slov.

Čisté slovo Božie.

Nič viac.

Sedem postáv, sedem hlasov.

Počujem otázku, ktorú kladie Hospodin: „Kto ti povedal, že si nahý? Nejedol si azda zo stromu, z~ktorého som ti zakázal jesť?“ Viem, že tá otázka patrí aj mne.

Mám zimomriavky z~davu, ktorý kričí: „Ukrižovať! Ukrižuj ho!!!“ A~pýtam sa sama seba, nestála by som v~tom dave vtedy i ja?

„Jeho krv na nás a na naše deti!“ ...strach mi zviera hrdlo.

Počujem slová výsmechu a pohŕdania: „Iných zachraňoval, sám seba nemôže zachrániť?“

Počujem aj hlas Ježiša: „Odpusť im, lebo nevedia, čo činia!“

V~tej riave lží a hnevu sa objaví záblesk pravdy a pochopenia od človeka, od ktorého by to nik nečakal: „Ani ty sa nebojíš Boha, hoci si takisto odsúdený? Lenže my spravodlivo dostávame zaslúžený trest, ale on neurobil nič zlé.“
A~Ježišove slová: „Otče, odpusť im, lebo nevedia, čo robia.“

Cítim Ježišov strach.
„Éli, Éli, lema sabachthani?“ -- „Bože môj, Bože môj, prečo si ma opustil?“

„Je dokonané!”

Nie je to príbeh o~nich. Je to príbeh o~nás.
\autor{Slávka Kráľová}


\clanok{Jarná víkendovka v~Berei}
17.~--~19.~4. sa konala zborová víkendovka v~Modre v~Berei s~témou „Spolu pracovať, žiť, uctievať“. Počasie bolo veľmi príjemné, program obohacujúci i oddychový, jedlo inovatívne. V~piatok bola opekačka a v~sobotu sme sa stretli (už v~hojnejšom počte) pri pestrých aktivitách.

Prvá téma bola doobeda a mal ju J.~Szőllős -- „Ako spolu pracovať“. Okrem iného sme si spoločne vypísali činnosti a služby, ktoré v~zbore máme, pravidelné aj občasné. Cieľom toho bolo a aj bude vytvorenie a prezentovanie katalógu služieb, ktorý môže pomôcť novým ľuďom aktívne sa zapojiť do zborového „spolupracujeme“.

Nasledovala spoločná prechádzka do krásnej prírody („Zámčisko“), počas ktorej bol dobrý priestor na rozhovory a spoznanie sa aj s~priateľmi zboru, ktorí využili možnosť prísť na víkendovku.

Poobede sme mohli hrať alebo sledovať športové aktivity (hlavne volejbalové zápasy) alebo zúčastniť sa výborne zostaveného biblického kvízu.

Deti mali radosť z~besiedky, zo vzájomných hier, tvorivých prác i z~venčenia psov.

Podvečer mal tému „Ako spolu žiť“ F.~Barkóczi.

Uzavretie témy mal D.~M.~Chuchút -- „Ako spolu uctievať“ -- v~nedeľu doobeda na Palisádach.
Sme vďační Pánu Bohu, že sme mohli mať ako spoločenstvo takýto priestor budovať sa a viac byť SPOLU.
\autor{Evka Syčová}


\clanok{Zbierky v~apríli}
\table{lr}{
Na misiu (Černihiv) & 942,00~€\cr
Na investičný fond	& 604,50~€\cr
}
\vskip1em

Pripomíname, že zbierky sa u~nás konajú pravidelne takto:
\vskip-1ex\begitems
* každú 2.~nedeľu v~mesiaci je zbierka venovaná misii a
* každú 4.~nedeľu je zbierka na investície.
\enditems

{\bf Pre dosporenie na rekonštrukciu} sme vypočítali príspevok na člena vo výške 150~€. Prispieť je možné formou zbierky počas 4.~nedieľ v~mesiaci alebo bankovým prevodom na zborový účet s~uvedením variabilného symbolu~777, prípadne aj osobne s.~Ľubke Kohútovej.

Bankové spojenie: SK36 0900 0000 0000 1147 1836, SWIFT: GIBASKBX

\n  1.	5.	Milica	MALÁ;
\n  1.	5.	Andrea	ČURILLOVÁ;
\n  3.	5.	Dárius	KRÁĽ;
\n  4.	5.	Peter	BUZÁŠ ml.;
\n  8.	5.	Vladimír	KRAJČÍ;
\n  9.	5.	Ján	ŠTEFKO;
\n 11.	5.	Želmíra	PRAŽENICOVÁ;
\n 16.	5.	Ján	SZŐLLŐS;
\n 17.	5.	Lenka	KOVÁČOVÁ;
\n 18.	5.	Anna	DANTEROVÁ;
\n 19.	5.	Oľga	VALCHÁŘOVÁ;
\n 20.	5.	Rastislav	PAULEN;
\narodeniny


\program{
\p  1 ; pi ; 17.30 ; Dorast ;.;;
\p  2 ; so ; 18.00 ; Mládež ;.;;
\p  3 ; ne ;  9.30 ; Bohoslužby (J.~Szőllős + VP) ;.;;
\p  4 ; po ; 18.00 ; Modlitebná skupinka / skupinka „Rastieme v~Kristovi“ ;.;;
\p  5 ; ut ; 15.00 ; Biblická hodina pre seniorov (P.~Pivka) ;.;;
\p  6 ; st ;.;;.;;
\p  7 ; št ; 18.00 ; Biblická hodina (D.~M.~Chuchút) ;.;;
\p  8 ; pi ; 17.30 ; Dorast ;.;;
\p  9 ; so ; 18.00 ; Mládež ;.;;
\p 10 ; ne ;  9.30 ; Bohoslužby (F.~Barkóczi) ; 12.00 ; Spoločný obed v~SED ;
\p 11 ; po ; 18.00 ; Modlitebná skupinka ;.;;
\p 12 ; ut ; 15.00 ; Biblická hodina pre seniorov (P.~Pivka) ;.;;
\p 13 ; st ;.;;.;;
\p 14 ; št ; 18.00 ; Biblická hodina (F.~Barkóczi) ;.;;
\p 15 ; pi ; 17.30 ; Dorast ;.;;
\p 16 ; so ; 18.00 ; Mládež ;.;;
\p 17 ; ne ;  9.30 ; Bohoslužby (V.~Pototskyi / D.~M.~Chuchút) ;.;;
\p 18 ; po ; 18.00 ; Modlitebná skupinka / skupinka „Rastieme v~Kristovi“ ;.;;
\p 19 ; ut ; 15.00 ; Biblická hodina pre seniorov (P.~Pivka) ;.;;
\p 20 ; st ; 17.30 ; Stretnutie sestier (B.~Dobová) ;.;;
\p 21 ; št ; 18.00 ; Biblická hodina (D.~M.~Chuchút) ;.;;
\p 22 ; pi ; 17.30 ; Dorast ;.;;
\p 23 ; so ; 18.00 ; Mládež ;.;;
\p 24 ; ne ;  9.30 ; Bohoslužby (D.~M.~Chuchút) ;.;;
\p 25 ; po ; 18.00 ; Modlitebná skupinka ;.;;
\p 26 ; ut ; 15.00 ; Biblická hodina pre seniorov (P.~Pivka) ;.;;
\p 27 ; st ;.;;.;;
\p 28 ; št ; 18.00 ; Biblická hodina (F.~Barkóczi) ;.;;
\p 29 ; pi ; 17.30 ; Dorast ;.;;
\p 30 ; so ; 18.00 ; Mládež ;.;;
\p 31 ; ne ;  9.30 ; Bohoslužby (bude oznámené neskôr) ;.;;
}

\tiraz
\bye
